\documentclass[12pt,a4paper]{article}
\usepackage[UTF8]{ctex}
\usepackage{float}
\usepackage{amsmath}
\usepackage{enumerate}
\usepackage{booktabs}
\usepackage{graphicx}
\usepackage{longtable}
\usepackage{subfigure}

\usepackage{url}

% for plotting 
\usepackage{caption}
\usepackage{pgfplots}

% for pseudo code 
\usepackage{algorithm}
\usepackage[noend]{algpseudocode}

% for reference 
\usepackage{hyperref}
\usepackage{cleveref}
\crefname{equation}{方程}{方程}
\Crefname{equation}{方程}{方程}
\crefname{table}{表}{表}
\Crefname{table}{表}{表}
\crefname{figure}{图}{图}
\Crefname{figure}{图}{图}

% for code 
\usepackage{listings}
\usepackage{xcolor}
\usepackage{inconsolata}
\definecolor{darkgreen}{rgb}{0,0.6,0}

\lstset {
    basicstyle=\footnotesize\ttfamily, % basic font setting
    breaklines=true, 
    frame=single,     % {single, shadowbox, bottomline}
    keywordstyle=\color{blue}, 
    commentstyle=\color{darkgreen},
    stringstyle=\color{red},
    showstringspaces=false,
    % backgroundcolor=\color{black!5}, % set backgroundcolor
    % numbers=left, 
    % numberstyle=\consolas,
}

% Microsoft Word A4 paper default layout 
\usepackage[a4paper, left=3.18cm, right=3.18cm, top=2.54cm, bottom=2.54cm]{geometry}

\captionsetup[figure]{name={图}}
\captionsetup[table]{name={表}}

\title{DIP Lab 2: Image Warping}
\author{2017011620 \quad 计73 \quad 李家昊}
\date{\today}

\begin{document}

\maketitle

\section{Perspective Warping}

首先展示生成效果，如\Cref{fig:perspective}所示。
\begin{figure}[H]
    \centering
    \subfigure[Source]{
        \includegraphics[width=0.4\textwidth]{../image/source.jpg}
    }
    \subfigure[Target]{
        \includegraphics[width=0.4\textwidth]{../image/target.jpg}
    }
    \subfigure[Output]{
        \includegraphics[width=0.8\textwidth]{../output/warp_perspective.jpg}
    }
    \caption{Perspective Warping}
    \label{fig:perspective}
\end{figure}

核心算法如下，首先需要得到变换矩阵，设其为，
\begin{equation}
    \boldsymbol{A} = \left(
    \begin{matrix}
        a & b & c\\
        d & e & f\\
        g & h & i\\
    \end{matrix}
    \right)
\end{equation}
在原图中找到四点的坐标为$\{(x_k, y_k)\}_{k=1}^4$，在目标图像中找到它们对应的四点的坐标$\{(u_k, v_k)\}_{k=1}^4$，则在齐次坐标下有，
\begin{equation}
    \left(
    \begin{matrix}
        a & b & c\\
        d & e & f\\
        g & h & i\\
    \end{matrix}
    \right)
    \left(
    \begin{matrix}
        x_1 & x_2 & x_3 & x_4\\
        y_1 & y_2 & y_3 & y_4\\
        1 & 1 & 1 & 1\\
    \end{matrix}
    \right)
    = \left(
    \begin{matrix}
        u_1 & u_2 & u_3 & u_4\\
        v_1 & v_2 & v_3 & v_4\\
        1 & 1 & 1 & 1\\
    \end{matrix}
    \right)
\end{equation}
其中两边取转置，即，
\begin{equation}
    \left(
    \begin{matrix}
        x_1 & y_1 & 1\\
        x_2 & y_2 & 1\\
        x_3 & y_3 & 1\\
        x_4 & y_4 & 1\\
    \end{matrix}
    \right)
    \left(
    \begin{matrix}
        a & b & c\\
        d & e & f\\
        g & h & i\\
    \end{matrix}
    \right)^T
    = \left(
    \begin{matrix}
        u_1 & v_1 & 1\\
        u_2 & v_2 & 1\\
        u_3 & v_3 & 1\\
        u_4 & v_4 & 1\\
    \end{matrix}
    \right)
\end{equation}
通过解线性方程组即可求出变换矩阵$\boldsymbol{A}$，这样就可以对图像进行变换了。首先求出其逆矩阵$\boldsymbol{A}^{-1}$，然后扫描目标图像的每一个像素位置$(r_d, c_d)$，根据下式，
\begin{equation}
    (r_s, c_s, 1)^T = s \cdot \boldsymbol{A}^{-1} (r_d, c_d, 1)^T
\end{equation}
求出原图中对应的像素位置$(r_s, c_s)$，其中$s$为齐次坐标的归一系数，得出坐标可能不是整数，则四舍五入到最近的整数格点$(\text{round}(r_s), \text{round}(c_s))$，即最邻近插值，然后直接将原图此坐标下的值拷贝到目标图的对应位置。核心代码如下，
\begin{lstlisting}[language=Python]
def get_perspective_transform(src_rect, dst_rect):
    assert src_rect.shape == (4, 2) and dst_rect.shape == (4, 2)
    src_rect = np.concatenate([src_rect.astype(np.float32), np.ones((4, 1), dtype=np.float32)], axis=-1)
    dst_rect = np.concatenate([dst_rect.astype(np.float32), np.ones((4, 1), dtype=np.float32)], axis=-1)
    matrix, _, _, _ = np.linalg.lstsq(src_rect, dst_rect)
    return matrix.T

def warp_perspective(src, matrix, dst_size):
    assert src.ndim == 2 or src.ndim == 3
    src_h, src_w = src.shape[:2]
    dst_w, dst_h = dst_size
    dst_shape = (dst_h, dst_w, src.shape[2]) if src.ndim == 3 else (dst_h, dst_w)
    dst = np.zeros(dst_shape, dtype=np.uint8)
    inv = np.linalg.inv(matrix)
    for dst_y in range(dst_h):
        for dst_x in range(dst_w):
            src_x, src_y, src_scale = np.matmul(inv, [dst_x, dst_y, 1])
            src_x = int(round(src_x / src_scale))
            src_y = int(round(src_y / src_scale))
            if 0 <= src_x < src_w and 0 <= src_y < src_h:
                dst[dst_y, dst_x] = src[src_y, src_x]
    return dst
\end{lstlisting}

\section{Sphere Warping}

首先展示结果，如\Cref{fig:sphere}所示。
\begin{figure}[H]
    \centering
    \subfigure[Source]{
        \includegraphics[height=5.5cm]{../image/warping.png}
    }
    \subfigure[Output]{
        \includegraphics[height=5.5cm]{../output/warp_sphere.jpg}
    }
    \caption{Sphere Warping}
    \label{fig:sphere}
\end{figure}

核心算法如下，将原图均匀铺在一个水晶球上面，剖面图如\Cref{fig:sphere_alg}所示，
\begin{figure}[H]
    \centering
    \includegraphics[width=0.5\textwidth]{fig/sphere.pdf}
    \caption{Sphere Warping Algorithm}
    \label{fig:sphere_alg}
\end{figure}
设原图的中心为$(r_{s0}, c_{s0})$，目标图的中心为$(r_{d0}, c_{d0})$，水晶球的半径为$l_{d0} = \max(r_{d0}, c_{d0}) / (\pi / 2)$，对于目标图的每个像素$(r_d, c_d)$，首先计算出它到图像中心的距离$l_d= \sqrt{(r_d - r_{d0})^2 + (c_d - c_{d0})^2}$，由几何关系可得，原图对应点距离原图得中心点的距离为，
\begin{equation}
    l_s = l_{d0} \cdot \phi = l_{d0} \cdot \arcsin\left(\frac{l_d}{l_{d0}}\right)
\end{equation}
由相似三角形关系，可知原图对应点$(r_s, c_s)$满足，
\begin{equation}
    (r_s - r_{s0}, c_s - c_{s0}) = \frac{l_s}{l_d} (r_d - r_{d0}, c_d - c_{d0})
\end{equation}
这样一来，就可以求出目标图像的每一点对应的原图坐标$(r_s, c_s)$了，同样，如果坐标不为整数，则四舍五入到邻近的整数格点。

核心代码如下，
\begin{lstlisting}[language=Python]
def warp_sphere(src, dst_size):
    assert src.ndim == 2 or src.ndim == 3
    src_h, src_w = src.shape[:2]
    dst_w, dst_h = dst_size
    dst_shape = (dst_h, dst_w, src.shape[2]) if src.ndim == 3 else (dst_h, dst_w)
    dst = np.zeros(dst_shape, dtype=np.uint8)

    max_dst_radius = min(dst_w, dst_h) // 2
    dst_center_y, dst_center_x = max_dst_radius, max_dst_radius
    src_center_y, src_center_x = src_h // 2, src_w // 2

    for dst_y in range(dst_h):
        for dst_x in range(dst_w):
            dst_radius = math.sqrt((dst_y - dst_center_y) ** 2 + (dst_x - dst_center_x) ** 2)
            if dst_radius > max_dst_radius:
                continue
            src_radius = max_dst_radius * math.asin(dst_radius / max_dst_radius)
            scale_factor = src_radius / dst_radius if dst_radius else 1
            src_y = int(round(scale_factor * (dst_y - dst_center_y) + src_center_y))
            src_x = int(round(scale_factor * (dst_x - dst_center_x) + src_center_x))
            if 0 <= src_y < src_h and 0 <= src_x < src_w:
                dst[dst_y, dst_x] = src[src_y, src_x]
    return dst
\end{lstlisting}

\section{Extra Task}

这里做了一个拓展实验，前面的球面形变产生了凸透镜效果，于是我在这里尝试了一下凹透镜效果，如\Cref{fig:sphere_inv}所示，
\begin{figure}[H]
    \centering
    \subfigure[Source]{
        \includegraphics[height=5.5cm]{../image/warping.png}
    }
    \subfigure[Output]{
        \includegraphics[width=5.5cm]{../output/warp_sphere_inv.jpg}
    }
    \caption{Inverse Sphere Warping}
    \label{fig:sphere_inv}
\end{figure}

原理是把原图均匀铺在一个锥面上，剖面图如\Cref{fig:sphere_inv_alg}所示，
\begin{figure}[H]
    \centering
    \includegraphics[width=0.5\textwidth]{fig/sphere_inv.pdf}
    \caption{Inverse Sphere Warping Algorithm}
    \label{fig:sphere_inv_alg}
\end{figure}

同样设目标图一点$(r_d, c_d)$到中心距离为$l_d$，则原图对应点到中心距离为，
\begin{equation}
    l_s = l_{d0} \cdot \phi = l_{d0} \cdot \arccos\left(\frac{l_{d0} - l_d}{l_{d0}}\right)
\end{equation}
其余操作同凸透镜效果。核心代码类似，不再赘述，可以参考附录的完整代码。

\section{Appendix: Codes}

附上了完整代码，这里导入的opencv仅仅用来读取和保存图片。
\lstinputlisting[language=Python]{../warp.py}

\end{document}
